\subsection*{Both fixed sources in the complete inner and outer regions}

\statement{B.113}{Proved}{The whole annular regions have their actual fixed-source centralizers, including the selected outer source in the original cover}
\label{stmt:B113}
Retain the actual surface $\mathscr Q$, its involution
$F(q,t)=(q,-t)$, and its free double cover $Z$ from B.106.
Put
\[
 V_< =\{q\in\mathbb C:0<|q|<1\},\qquad
 V_> =\{q\in\mathbb C:|q|>1\},
\]
and write $\mathscr Q_<,\mathscr Q_>,Z_<,Z_>$ for their actual
inverse images. Every critical fiber in the specified region is
included. For $\epsilon\in\{<,>\}$ and $j\in\{0,\infty\}$,
write $D_j^\epsilon=D_j\cap\mathscr Q_\epsilon$. Use a base
point on the indicated fixed source and the resulting literal
based involution. Then its source inclusion is injective, induces
the full profinite topology, and satisfies
\begin{equation}\label{eq:B113-all-regional-sources}
 \begin{aligned}
 \pi_1(\mathscr Q_\epsilon)^F
       &=\pi_1(D_j^\epsilon),\\
 \bigl(\widehat{\pi_1(\mathscr Q_\epsilon)}\bigr)^F
       &=\overline{\pi_1(D_j^\epsilon)}
         \cong\widehat{\pi_1(D_j^\epsilon)}.
 \end{aligned}
\end{equation}
The ranks, with the base point chosen separately for each action,
are
\[
\begin{array}{c|cc}
 &D_0&D_\infty\\ \hline
 V_<&1&5\\
 V_>&5&1
\end{array}
\]
where rank one means the infinite cyclic group. We identify each
source group with its actual based image in
\eqref{eq:B113-all-regional-sources}; an unspecified change between
the two different fixed base points is not part of that notation.

For the original cover, take the selected outer base point
\[
 P_>=(1+i\sqrt3,1-i\sqrt3)\in S\cap Z_>,
 \qquad \bar P_>=(4,0)\in D_0^>.
\]
Its source $S_>=S\cap Z_>$ is connected and is an unramified
double cover of $D_0^>$. Its actual source map is injective with
its full profinite topology, and
\begin{equation}\label{eq:B113-original-outer}
 \pi_1(Z_>,P_>)^F=\pi_1(S_>,P_>)\cong F_9,
 \qquad
 \bigl(\widehat{\pi_1(Z_>,P_>)}\bigr)^F
       =\overline{\pi_1(S_>,P_>)}\cong\widehat F_9.
\end{equation}
These statements concern whole individual regions, with their old
annular loops retained. They do not identify the group obtained
by joining the regions through a neck, or the global fixed source.

\begin{proof}
\emph{B.113.1. The other cone in a retracted amalgam.}
We first prove the comparison needed to transfer B.112 from its
bad region to the whole region. Let $R$ be a discrete group with
an orientation character $o:R\to C_2$, and let $c\in R$ have
order two and $o(c)\ne1$. Suppose a specified subgroup $C\le R$
lies in $\ker o$ and has a retraction $r:R\to C$. Let
\[
 E=(C\times\langle d\rangle)*_C R,\qquad d^2=1,
\]
with the displayed edge inclusions. Extend orientation to $E$ by
sending $d$ to the nontrivial element of $C_2$. Then
\begin{equation}\label{eq:B113-other-vertex}
 C_E(c)=C_R(c),\qquad
 C_{\widehat E}(c)=C_{\widehat R}(c),
\end{equation}
where $\widehat R\to\widehat E$ is injective.

Indeed, the map $E\to R$ that is the identity on $R$ and kills
$d$ is a retraction. It proves the asserted full induced topology
and injectivity on completions. In the ordinary amalgam tree,
$c$ is odd and every conjugate edge subgroup is even. Thus $c$
fixes no edge. Since it fixes the $R$ vertex, its fixed subtree
is precisely that vertex. A centralizing element preserves it,
proving the first equality in \eqref{eq:B113-other-vertex}.

For the completed equality we give the cofinal comparison
explicitly. Regard $r$ as an idempotent endomorphism of $R$ using
the specified inclusion of $C$. Given a finite-index normal
$M\triangleleft R$, put
\[
 M_0=M\cap\ker o,\qquad N=M_0\cap r^{-1}(M_0).
\]
Then $N\subset M_0$ is finite-index normal and $r(N)\subset N$.
Set $A_N=R/N$ and let $D_N$ be the image of $C$ in $A_N$.
The induced retraction has image $D_N$, and the finite-vertex
amalgam
\[
 E_N=(D_N\times C_2)*_{D_N}A_N
\]
is a quotient of $E$. The image $c_N$ of $c$ remains an odd
involution, while its edge subgroup is even. The ordinary tree
therefore gives
\[
 C_{E_N}(c_N)=C_{A_N}(c_N).
\]
The right side is finite.

The regular-permutation construction in B.108.3 supplies a normal
finite-index free subgroup of finite rank in every $E_N$, and
proves its residual finiteness. The all-discrete-module comparison
and finite virtual $2$-cohomological type are exactly those verified
in B.108.4: a finite-rank free kernel, finite quotient, and two
nonzero kernel rows in the cohomology spectral sequence. We thus
apply only the prime-order \emph{element} conclusion of
\citet[Theorem~4.3, Sections~4.1--4.2]{BoggiZalesskii2025}, with
all its hypotheses, to obtain
\[
 C_{\widehat E_N}(c_N)=C_{A_N}(c_N)\subset A_N.
\]
No conclusion for a nonsolvable finite subgroup is used.

Every finite quotient $\rho:E\to Q$ factors through one of these
$E_N$. Apply the construction with
$M=\ker(\rho|_R)$. Then $N\subset M$, and the restriction to
$C\times C_2$ factors through $D_N\times C_2$, because
$N\cap C\subset\ker(\rho|_C)$. This proves cofinality for
arbitrary finite quotients of the whole amalgam.

Let $g\in C_{\widehat E}(c)$. Its image in every
$\widehat E_N$ lies in $A_N$, the image of $R$. If $g$ were
outside the compact image of $\widehat R$, a finite quotient of
$\widehat E$ would separate $g$ from that image. Factoring this
quotient through the corresponding $E_N$ contradicts the preceding
membership. Hence $g\in\widehat R$ and centralizes $c$ there.
This proves the second equality in \eqref{eq:B113-other-vertex}.
We have not assumed residual finiteness of the original $R$.

\emph{B.113.2. Application to the complete actual regions.}
B.107 constructs the actual collar decomposition over $V_<$,
including all critical fibers, and proves
\[
 E_<=(\mathbb Z\times C_{2,0})*_{\mathbb Z}R_<.
\]
Here $R_<$ is exactly the bad annular orbifold group of B.112;
the edge is the specified boundary section, its retraction is
projection to the base, and its image is ordinary. At a base
point on $D_\infty^<$ the cone $c$ is its actual odd involution.
The basepoint connector used in the collar gives the based form
of this same amalgam. Applying \eqref{eq:B113-other-vertex} and
then B.112 gives
\[
 C_{E_<}(c)=\pi_1(D_\infty^<)\times\langle c\rangle,
 \qquad
 C_{\widehat E_<}(c)
    =\widehat{\pi_1(D_\infty^<)}\times\langle c\rangle.
\]
The source carries its full topology in $R_<$ by B.112, and
$R_<$ carries its full topology in $E_<$ by the displayed
retraction. Intersecting with the ordinary index-two subgroup,
and with its open full completion, proves
\eqref{eq:B113-all-regional-sources} for $D_\infty^<$.
The $D_0^<$ case was proved in B.107.

The actual homeomorphism
\[
 J(q,t)=(1/\overline q,1/\overline t)
\]
commutes with $F$, exchanges $V_<$ and $V_>$, and exchanges
$0$ and $\infty$ in the fiber. It preserves every deleted
family: in particular
$1/\overline{T_u(q)}=T_u(1/\overline q)$ for the same $u$,
and it preserves the equation $qt^2=1$. It therefore transports
the two proved comparisons in $\mathscr Q_<$ to the indicated
comparisons in $\mathscr Q_>$. This is a comparison of the
quotient surfaces $\mathscr Q$; no lift of $J$ to a specified
involution on $Z$ is assumed. Base points may be transported
within the indicated fixed source by a chosen path in that
source. Such transport intertwines the literal based involutions,
so the equalities and full topology hold at every stated source
base point.

Finally, B.106 locates the four $\beta_u$ in $V_<$ and the four
$\alpha_u$ in $V_>$, all distinct. Thus $D_0^<$ and
$D_\infty^>$ are annuli, whereas $D_\infty^<$ and $D_0^>$
are annuli with four additional points removed. Their groups have
the four ranks displayed in the statement.

\emph{B.113.3. The selected original outer cover.}
We now use the original covering map itself. By B.106 its
restriction over $D_0$ is obtained by adjoining $\sqrt{1-q}$,
with ordered roots
\[
 v=1+\sqrt{1-q},\qquad w=1-\sqrt{1-q}.
\]
The point $P_>$ is in the actual surface: $q=4$ is not an old
base puncture or an $\alpha_u$, since $|\alpha_u|<3$, and the
two roots are distinct, nonzero, and different from $1$.
The loop $q=4e^{2\pi i s}$, $0\le s\le1$, lies in $D_0^>$.
Its image $1-q$ winds once around zero, so the square root
changes sign. The restricted two-sheeted cover is therefore
connected. Its source group is an index-two subgroup of $F_5$,
hence is free of rank $2(5-1)+1=9$. The same connected source
cover meets both sheets of $Z_>\to\mathscr Q_>$, so the whole
restricted cover is connected as well.

Let $\Pi=\pi_1(\mathscr Q_>,\bar P_>)$ and let
$\Pi'\le\Pi$ be the index-two subgroup supplied by $Z_>$.
Let $I=\pi_1(D_0^>,\bar P_>)\le\Pi$ and
$I'=\pi_1(S_>,P_>)$. The based pullback cover gives
\[
 I\cap\Pi'=I',\qquad [I:I']=2.
\]
Both actions are the actual involutions induced by $F$, since
the chosen base points are fixed and the covering is equivariant.
Intersecting the already proved equality $\Pi^F=I$ with
$\Pi'$ gives $(\Pi')^F=I'$.

Finite-index cores identify $\widehat{\Pi'}$ with the open
kernel of the covering character in $\widehat\Pi$. The source
$I$ has its full topology in $\Pi$ by the preceding comparison;
the same finite-index argument identifies $\widehat{I'}$ with
the corresponding open subgroup of $\widehat I$. Consequently
\[
 \overline I\cap\widehat{\Pi'}=\overline{I'},
 \qquad
 \overline{I'}\cong\widehat{I'}\cong\widehat F_9.
\]
Intersecting $(\widehat\Pi)^F=\overline I$ with the open
subgroup $\widehat{\Pi'}$ proves the second equality in
\eqref{eq:B113-original-outer}, including the required direction
of transfer and the source's full induced topology.
\end{proof}
